\clearpage
\part{第一题：基础实现}
本篇以C语言基础、枚举、统计和小模拟为主。先读懂对象、规则与边界，再把它们翻译成变量、数组和循环；基础实现能够清楚表达规则时，优先使用基础实现。
\section{程序基本功与做题技巧}
\subsection{程序基本功}
程序基本功就是C语言的基础知识，包括格式化输入输出、数据的存储与运算（整型、浮点数）、函数的编写与调用、数组、结构体的使用等。下面按做题中实际会用到的顺序简单总结。
\subsubsection{格式化输入输出}
scanf按给定格式读入变量，printf按给定格式输出结果。读入普通变量时通常要写地址符\&；字符数组的名字已经可以表示首元素地址。格式符必须匹配变量类型。
\begin{longtable}{P{2.3cm}P{3.7cm}P{4.5cm}P{3.7cm}}
\toprule 类型 & 常见用途 & scanf格式 & printf格式\\\midrule\endhead
int & 下标、计数、普通整数 & \texttt{\%d} & \texttt{\%d}\\
long long & 大整数、前缀和、乘积 & \texttt{\%lld} & \texttt{\%lld}\\
double & 实数、均值、方差 & \texttt{\%lf} & \texttt{\%f}，如\texttt{\%.6f}\\
char & 单个字符、类别 & \texttt{\%c} & \texttt{\%c}\\
char数组 & 单词或一整行文本 & \texttt{\%100s}或fgets & \texttt{\%s}\\
\bottomrule\end{longtable}
\begin{lstlisting}
int n; long long total; double x; char type;
scanf("%d%lld%lf %c", &n, &total, &x, &type);
printf("%d %lld %.6f %c\n", n, total, x, type);
char word[101];
scanf("%100s", word);
\end{lstlisting}
\texttt{\%c}默认不跳过空白；格式串中在它前面加空格，表示先跳过空格与换行。\texttt{\%s}遇到空白停止，因此不能读取含空格的整行。fgets读取整行时可能保留换行；数字输入后再读整行，要处理该行剩余字符。输出空格和换行按题面要求写，不输出“请输入”等提示语。
\subsubsection{整型、浮点数与运算}
本讲义所用评测环境中int通常占4字节，范围为$-2147483648$到$2147483647$；long long占8字节，范围约为$\pm9.22\times10^{18}$。不要只看最后答案，\textbf{每个中间结果也必须装得下}。\texttt{long long s = a*b;}仍可能先进行int乘法，应写\texttt{1LL*a*b}。

整数除法舍去小数部分（向0截断），如\texttt{5/2}为2；需要实数结果时写\texttt{5.0/2}或\texttt{(double)a/b}。\texttt{\%}用于整数取余，不能用于double。对非负整数x，向上取整$x/d$可用\texttt{x/d + (x\%d != 0)}，避免\texttt{x+d-1}溢出。double有舍入误差，通常不直接用\texttt{==}比较计算所得的实数；误差判定方式要与题目精度要求一致。比较均值与整数阈值时，若可以安全交叉相乘，优先使用整数比较。
\begin{lstlisting}
int a = 1000000000, b = 3;
long long product = 1LL * a * b;
double mean = (double)a / b;
\end{lstlisting}
\subsubsection{函数的编写与调用}
把一个明确任务写成函数：参数说明输入，返回值说明结果。先定义函数再调用，或在前面声明。C语言参数按值传递，修改函数内的普通参数不会自动修改调用方；需要修改调用方变量时传指针。数组作为参数时，函数收到首元素地址，必须另传长度。
\begin{lstlisting}
long long sum_array(const int a[], int n) {
    long long s = 0;
    for (int i = 0; i < n; i++) s += a[i];
    return s;
}
void swap_int(int *x, int *y) {
    int t = *x; *x = *y; *y = t;
}
// Call: total = sum_array(a, n); swap_int(&a[0], &a[1]);
\end{lstlisting}
函数名应说明职责，例如判断、计数、更新；不要让同一个函数既读入又排序又输出所有状态。递归还必须有终止条件，搜索时另需撤销分支修改。
\subsubsection{一维数组、二维数组与初始化}
长度为n的数组用0下标时合法范围是$0$到$n-1$；用1下标时至少申请$n+1$格。字符数组还需要末尾的\texttt{'\textbackslash0'}。二维数组\texttt{a[i][j]}通常让i表示行、j表示列，循环边界分别对应行数和列数。
\begin{lstlisting}
const int N = 1005;
int a[N], grid[N][N]; // Global arrays are initially zero.
int count[10] = {0};
for (int i = 0; i < n; i++) scanf("%d", &a[i]);
\end{lstlisting}
全局数组自动初始化为0，局部普通数组未初始化时不能直接读取。大数组宜放全局，避免局部栈空间不足。最小值、最大值通常用第一个实际元素初始化；若全为正数，把最小值初始化成0会出错。频率数组只能在值域可控时按值索引，负数需要偏移，巨大值域不能直接开数组。
\subsubsection{结构体与对象属性}
一个对象有多个属性时，用结构体把它们放在一起。例如点有横坐标、纵坐标和类别，矩形有四个边界。多个对象用结构体数组存储，读题时的名词便有了明确的数据表示。
\begin{lstlisting}
struct Point { int x, y; char type; };
struct Point p[1005];
scanf("%d%d %c", &p[0].x, &p[0].y, &p[0].type);
\end{lstlisting}
\subsection{读题重点}
\subsubsection{观察数据范围：两个首要检查}
\textbf{第一，数组不能开小。}若$n\le1000$且使用1下标，数组不能只有1000格；二维边界、可能访问的$n+1$位置、字符串末尾都要留出来。数组大小由\textbf{索引最大值}决定，不一定由输入个数决定：即使只输入500个编号，若编号可达$10^9$，也不能直接用这么大的编号作为普通数组下标。

\textbf{第二，运算过程中会不会超过int或long long的范围。}若$n=2\times10^5$且每项绝对值不超过$10^9$，总和上界为$2\times10^{14}$，需要long long；两项乘积的上界应相乘估算。即便long long也可能溢出，例如$10^{18}\times5$还安全，而$10^{18}\times100$已经越界。
\subsubsection{留意时间限制：时间复杂度与估算}
时间复杂度描述操作次数随输入规模增长的数量级，不是程序运行的精确秒数。估算时先看循环执行多少次，再忽略常数倍和低阶项。以下各段代码独立示意，假设n、m已给定，数组a、grid已读入，各次加法、比较的代价近似相同。

\Needspace{9\baselineskip}\textbf{一次遍历：$O(n)$。}
\begin{lstlisting}
long long sum = 0;
for (int i = 0; i < n; i++) {
    sum += a[i];
}
\end{lstlisting}
循环执行n次，每次做一次加法。n扩大到原来的10倍，主要操作次数也约为原来的10倍。

\Needspace{9\baselineskip}\textbf{前后两次遍历：仍是$O(n)$。}
\begin{lstlisting}
long long sum = 0;
for (int i = 0; i < n; i++) sum += a[i];
int positive = 0;
for (int i = 0; i < n; i++) {
    if (a[i] > 0) positive++;
}
\end{lstlisting}
两个循环依次执行，总共约2n次，忽略常数2后是$O(n)$。不能看到两个for就判断为平方复杂度。

\Needspace{9\baselineskip}\textbf{嵌套循环：$O(n^2)$。}
\begin{lstlisting}
long long equalPairs = 0;
for (int i = 0; i < n; i++) {
    for (int j = i + 1; j < n; j++) {
        if (a[i] == a[j]) equalPairs++;
    }
}
\end{lstlisting}
内层依次执行$n-1,n-2,\ldots,0$次，合计$n(n-1)/2$次比较，数量级为$O(n^2)$。n扩大10倍，比较次数约扩大100倍；内层不必每次都执行n次，仍可能是平方复杂度。

\Needspace{9\baselineskip}\textbf{两层循环长度不同：$O(nm)$。}
\begin{lstlisting}
long long sum = 0;
for (int i = 0; i < n; i++) {
    for (int j = 0; j < m; j++) {
        sum += grid[i][j];
    }
}
\end{lstlisting}
外层n次，每次内层m次，共nm次加法。若m始终是固定的10，便是10n次，可以记作$O(n)$；若m也随输入变化，就保留$O(nm)$。

\Needspace{9\baselineskip}\textbf{每次缩小一半：$O(\log n)$。}
\begin{lstlisting}
int steps = 0;
for (int length = n; length > 1; length /= 2) {
    steps++;
}
\end{lstlisting}
例如n=16时，length依次为16、8、4、2，循环4次，因为$2^4=16$。一般约执行$\log_2 n$次；n翻倍只会多执行约一次。

判断能否在时限内完成，要把\textbf{最大数据规模}代入。例如$n=10^5$时，一次遍历约$10^5$次，枚举所有数对约$5\times10^9$次，后者通常无法在几秒内完成。操作类型、机器速度和输入输出都会影响耗时，因此次数估算用于筛选算法，不能当作精确的用时保证。
\subsubsection{留意空间限制：空间复杂度与估算}
空间复杂度描述存储需求随规模增长的数量级。本书说明程序使用的主要数组，并把输入存储也计入；额外空间另有需要时会单独指出。一个n项数组为$O(n)$，$n\times m$矩阵为$O(nm)$，只保留若干计数变量为$O(1)$。递归栈也占空间，深度n的递归可能占$O(n)$栈空间。

实际估算为“元素个数$\times$每个元素字节数”，再把所有数组、容器和递归开销相加。$1000\times1000$的int矩阵约$4\times10^6$字节（约3.8MiB），两个矩阵约7.6MiB；$10^5$项long long数组约0.76MiB；$10^5\times10^5$的int矩阵约37GiB，显然不可行。结构体可能有对齐填充，STL容器还有管理开销，按实际\texttt{sizeof}和容量估算更准确。
\subsubsection{写完后的检查顺序}
先查格式与类型，再查最小规模、最大下标、重复值、负值、端点和更新顺序；最后分析算法为什么不重不漏。样例通过不代表所有数据正确。小规模时可写直接枚举程序作为对照。
